\documentclass{ltjsbook}
\usepackage{graphicx}
\usepackage{enumerate}
\usepackage{makeidx}
\makeindex
\title{索引のテスト}
\author{山田 太郎}
\date{\today}
\begin{document}
\maketitle
\chapter{\LaTeX の索引機能}
教科書や参考書のような書籍を読んでいると、索引\index{さくいん@索引}を見かけることが多いでしょう。索引が付いていると、読者が重要単語の書かれたページだけを簡単に拾い読みできるようになり非常に便利です。\LaTeX \index{らてふ@\LaTeX} では、索引を自動生成してくれるとても便利な機能があります。

\section{事前設定}
索引を付けたい文書では、以下の設定をする必要があります。
\begin{enumerate}[(1)]
\item
\verb|\usepackage{makeidx}| で、 \verb|makeidx.sty| を読み込む。
\item
プリアンブル部\index{ぷりあんぶるぶ@プリアンブル部}に \verb|\makeindex| を記載する。
\item
本文中の索引に収録したい言葉の直後に \verb|\index{よみ@索引項目}| を書いてマーク\index{まーく@マーク}。
- 索引を出力する位置 (通常は文末、つまり \verb|\end{document}| の直前) に \verb|\printindex| を記入。
\end{enumerate}

そして、最終的に文書を出力すると、そこには索引ページが自動生成されます。

\section{索引があると嬉しいこと}
なんといっても、索引があると、読者が重要だと思う単語のページだけを拾い読み\index{ひろいよみ@拾い読み}しやすくなり、とても便利なのです。
Microsoft Office\index{まいくろそふとおふぃす@Microsoft Office}のWord\index{わーど@Word}でも索引作成機能\index{さくいんさくせいきのう@索引作成機能}がありますが、\LaTeX を使って作った索引は、なんといってもレイアウト\index{れいあうと@レイアウト}などが完全自動調整されるので、まるで教科書のような美しい見た目になります。

\section{単語の収録ページ数と順番は自動調整される}
そして、索引に収録したページ数は自動的にふられ、収録順はあいうえお順 (辞書式順序\index{じしょしきじゅんじょ@辞書式順序}) に自動的に並び替えられます。手入力\index{てにゅうりょく@手入力}で索引を作ろうと思うと、それはかなりしんどい作業となりますが、やはりそこは\LaTeX の機能におまかせしてしまおうという考え方です。
\printindex
\end{document}

